\begin{frame}{Introduction}\framesubtitle{Contexte}
  Termith (Terminologie et indexation de textes en sciences humaines)~:
  \begin{itemize}
    \item{Projet \textsc{Anr}}
    \item{Multi-partenaires~:}
      \begin{figure}
        \centering
        \begin{tikzpicture}[node distance=2em]
          \node [] (inist) {\includegraphics[scale=.33]{logo/inist_min.eps}};
          \node [above=of inist.north west, xshift=-2em] (atilf) {\includegraphics[scale=.33]{logo/atilf_min.eps}};
          \node [above=of inist.north east, xshift=2em] (lidilem) {\includegraphics[scale=.5]{logo/lidilem_min.eps}};
          \node [below=of inist.south west, xshift=-2em] (lina) {\includegraphics[scale=.33]{logo/lina_min.eps}};

          \node [below=of inist.south east, xshift=2em] (inria) {\includegraphics[scale=.33]{logo/inria_min.eps}};

          %\node [draw=YellowOrange, fill=Goldenrod, star, star points=5, above=of atilf.north west, xshift=.25em, yshift=-2.5em, scale=.66] (leader) {};

          \path[->] (atilf) edge (inist);
          \path[->] (lina) edge (inist);
          \path[->] (lidilem) edge (inist);
          \path[->] (inria) edge (inist);
        \end{tikzpicture}
      \end{figure}
    \item{Centré sur l'activité de l'Inist}
  \end{itemize}
\end{frame}

\begin{frame}{Introduction}\framesubtitle{Contexte}
  \begin{block}<+->{Inist (Institut de l'information scientifique et technique)}
    \begin{itemize}
      \item{}
      \item{}
      \item{}
    \end{itemize}
  \end{block}

  \begin{block}<+->{Base de données bibliographiques}
    Test
  \end{block}
\end{frame}

